%==============================================================================
%  main.tex — 심리음향 모델 (Psychoacoustic Model) 기술 문서
%------------------------------------------------------------------------------
%  권장 컴파일 엔진 : XeLaTeX (기본) 또는 LuaLaTeX
%  한국어 조판       : kotex 패키지 + fontspec (CJK 폰트 필요)
%  권장 폰트         : Noto Serif CJK KR / Noto Sans CJK KR 또는 나눔(Nanum) 계열
%
%  컴파일 방법 (상호참조·목차 갱신을 위해 두 번 실행):
%      xelatex main.tex
%      xelatex main.tex
%  또는:
%      lualatex main.tex   (두 번)
%      latexmk -xelatex main.tex
%
%  폰트가 없거나 다른 폰트를 쓰려면 아래 preamble 의
%  \setmainhangulfont / \setmainfont 주석을 참고하여 교체하십시오.
%
%  주의: 이 문서를 작성한 샌드박스에는 TeX 배포판과 CJK 폰트가 설치되어 있지
%  않으므로 실제 PDF 컴파일은 외부 환경에서 수행해야 합니다. 자세한 내용은
%  README.md 를 참고하십시오.
%==============================================================================
\documentclass[11pt,a4paper]{article}

%--- 한국어 조판 -------------------------------------------------------------
% kotex 는 XeLaTeX/LuaLaTeX 에서 fontspec 을 통해 한글 폰트를 설정합니다.
\usepackage{kotex}

%--- 페이지 형상 -------------------------------------------------------------
\usepackage[margin=2.5cm]{geometry}

%--- 수학 -------------------------------------------------------------------
\usepackage{amsmath}
\usepackage{amssymb}
\usepackage{mathtools}
\usepackage{siunitx}
\sisetup{detect-all=true}
% siunitx v3 는 bit/byte 단위를 기본 제공하지 않으므로 명시적으로 선언한다.
\DeclareSIUnit\bit{bit}
\DeclareSIUnit\byte{B}

%--- 표 ---------------------------------------------------------------------
\usepackage{booktabs}
\usepackage{array}

%--- 그림 -------------------------------------------------------------------
\usepackage{graphicx}
\usepackage{tikz}
\usepackage{pgfplots}
\pgfplotsset{compat=1.18}

%--- 알고리즘 / 의사코드 ------------------------------------------------------
\usepackage{algorithm}
\usepackage{algpseudocodex}

%--- 하이퍼링크 및 상호참조 ---------------------------------------------------
\usepackage[hidelinks]{hyperref}
\usepackage{cleveref}

%--- 폰트 설정 (필요 시 주석 해제하여 교체) ------------------------------------
% 기본적으로 kotex 는 사용 가능한 한글 폰트를 자동 선택합니다.
% 특정 폰트를 강제하려면 아래 예시처럼 지정하십시오. (폰트가 설치되어 있어야 함)
%
%   \setmainfont{Noto Serif}            % 라틴 문자 본문 폰트
%   \setmainhangulfont{Noto Serif CJK KR} % 한글 본문 폰트
%   \setsanshangulfont{Noto Sans CJK KR}
%   % 나눔 계열 대안:
%   % \setmainhangulfont{NanumMyeongjo}
%   % \setsanshangulfont{NanumGothic}

%==============================================================================
%  재사용 매크로 (반복되는 표기법)
%==============================================================================
\newcommand{\dB}{\ensuremath{\,\mathrm{dB}}}          % 데시벨
\newcommand{\dBSPL}{\ensuremath{\,\mathrm{dB\,SPL}}}  % dB SPL
\newcommand{\Bark}{\ensuremath{\,\mathrm{Bark}}}      % Bark 단위
\newcommand{\Hz}{\ensuremath{\,\mathrm{Hz}}}          % Hz
\newcommand{\SPL}{\mathrm{SPL}}                        % 음압 레벨
\newcommand{\ATH}{T_{q}}                               % 절대 청취 문턱값
\newcommand{\SMR}{\mathrm{SMR}}                        % signal-to-mask ratio
\newcommand{\MNR}{\mathrm{MNR}}                        % mask-to-noise ratio
\newcommand{\SNR}{\mathrm{SNR}}                        % signal-to-noise ratio
\newcommand{\dz}{\Delta z}                             % Bark 축 거리
\DeclareMathOperator{\arctanop}{arctan}

%==============================================================================
\title{\textbf{심리음향 모델(Psychoacoustic Model)}\\[0.3em]
       \large 배경·수식·알고리즘 및 구현 기법}
\author{기술 조사 문서}
\date{\today}

\begin{document}
\maketitle

%------------------------------------------------------------------------------
\begin{abstract}
\noindent
본 문서는 지각적 오디오 부호화(perceptual audio coding)의 핵심 이론인 \emph{심리음향
모델}을 정리한다. 먼저 인간 청각의 생리적 특성과 청각 마스킹(auditory masking)에서
출발하여 심리음향 모델이 왜 필요한지(\Cref{sec:background})를 설명한다. 이어서 절대
청취 문턱값(Terhardt), Bark(임계대역) 척도 변환(Zwicker), 확산 함수(spreading
function), 마스킹 문턱값의 합성, 그리고 신호 대 마스크 비(SMR)/마스크 대 잡음
비(MNR) 등 표준 수식(\Cref{sec:math})을 유도·정리한다. 마지막으로 ISO/IEC 11172-3
(MPEG-1 Audio)의 심리음향 모델 1과 모델 2를 알고리즘 형태(\Cref{sec:algorithm})로
제시하고, 비트 할당(bit allocation)과의 결합 방식을 논의한다. 모든 수식의 기호는
정의와 함께 제시하며, 인용은 실재하는 표준·문헌만을 사용한다.
\end{abstract}

\tableofcontents
\bigskip

%==============================================================================
\section{배경 및 동기}\label{sec:background}
%==============================================================================

\subsection{왜 심리음향 모델이 필요한가}
디지털 오디오를 무손실로 저장하려면 매우 큰 비트율이 필요하다. 예컨대 CD 품질
(\SI{44.1}{\kilo\hertz} 표본화, 16비트, 스테레오)은 약 \SI{1.41}{\mega\bit\per\second}
를 요구한다. 손실 압축(lossy compression)은 \emph{사람이 실제로 지각하지 못하는}
성분을 제거하여 비트율을 크게 낮춘다. 무엇을 지각하지 못하는지를 정량적으로
판정하는 도구가 바로 심리음향 모델이며, 이는 \emph{지각적 무의미성}(perceptual
irrelevance)을 활용한다. 즉 양자화 잡음(quantization noise)을 사람이 들을 수 없는
수준, 곧 \emph{마스킹 문턱값}(masking threshold) 아래로 유지하면, 원신호와 지각적으로
구별되지 않으면서도 데이터량을 획기적으로 줄일 수 있다.

\subsection{인간 청각의 기초}
소리는 외이(outer ear)와 중이(middle ear)를 거쳐 내이(inner ear)의
달팽이관(cochlea, 와우)으로 전달된다. 달팽이관 내부의 \emph{기저막}(basilar
membrane)은 위치에 따라 서로 다른 주파수에 공명한다. 기저막의 바닥(base) 부근은
고주파에, 꼭대기(apex) 부근은 저주파에 최대로 반응한다. 이러한 \emph{장소
부호화}(tonotopic organization) 때문에 청각계는 일종의 주파수 분석기로
동작한다.

기저막의 각 지점은 유한한 대역폭을 갖는 대역통과 필터처럼 거동하며, 이를
\emph{임계대역}(critical band)이라 한다. 임계대역은 인간 청각의 주파수 분해능을
반영하는 비선형 척도이며, 이를 균등한 간격으로 만든 것이 Bark 척도이다. 한 임계대역
안에 놓인 성분들은 서로 강하게 상호작용(특히 마스킹)한다.

\subsection{절대 청취 문턱값(ATH)}
\emph{절대 청취 문턱값}(Absolute Threshold of Hearing, ATH) 혹은 정적 문턱값(threshold
in quiet)은 조용한 환경에서 특정 주파수의 순음이 들리기 시작하는 최소 음압
레벨이다. ATH는 주파수에 강하게 의존하여 \SIrange{2}{5}{\kilo\hertz} 부근에서 가장
낮고(가장 민감하고) 저주파·고주파로 갈수록 급격히 상승한다. 부호화기는 ATH 아래의
성분을 아예 부호화하지 않아도 무방하다.

\subsection{청각 마스킹}
\emph{마스킹}(masking)은 한 소리(마스커, masker)의 존재 때문에 다른 소리(피마스커,
maskee)의 청취 문턱값이 상승하는 현상이다. 두 가지로 구분된다.
\begin{itemize}
  \item \textbf{동시 마스킹}(simultaneous / frequency masking): 마스커와 피마스커가
        동시에 존재할 때, 마스커 주변 주파수의 문턱값이 상승한다. 마스킹 곡선은
        Bark 축에서 대략 삼각형 모양이며 저주파 방향보다 고주파 방향으로 더 넓게
        퍼진다(비대칭).
  \item \textbf{시간적 마스킹}(temporal masking): 마스커가 사라진 뒤에도 짧은 시간
        (약 \SIrange{50}{200}{\milli\second}) 동안 문턱값이 유지되는 후행
        마스킹(post-masking)과, 마스커 시작 직전 짧은 구간(약
        \SI{5}{\milli\second})에서 나타나는 선행 마스킹(pre-masking)이 있다.
\end{itemize}
마스커가 순음에 가까운 톤성(tonal)인지 잡음에 가까운 비톤성(non-tonal)인지에 따라
마스킹 능력이 다르며, 이 차이는 \Cref{sec:math,sec:algorithm}에서 정량화한다.

\subsection{지각적 부호화에서의 역할}
MPEG-1 Layer III(MP3), MPEG-2/4 AAC 등 지각적 오디오 부호화기는 심리음향 모델로부터
각 주파수 대역(또는 서브밴드)의 \emph{마스킹 문턱값}을 얻는다. 그런 다음 그 문턱값
아래로 양자화 잡음을 유지하도록 비트를 할당한다. 결과적으로 지각적으로 중요한
대역에는 비트를 많이, 마스킹되어 들리지 않는 대역에는 적게 배분함으로써 동일 비트율
에서 최상의 지각 품질을 얻는다.

%==============================================================================
\section{수식}\label{sec:math}
%==============================================================================
이 절에서는 심리음향 모델에 사용되는 표준 수식을 정의와 함께 제시한다. 별도 언급이
없으면 주파수 $f$의 단위는 \si{\hertz}, 레벨의 단위는 \si{\decibel}(음압 레벨의 경우
\dBSPL)이다.

\subsection{절대 청취 문턱값 (Terhardt)}
주파수 $f$에서의 절대 청취 문턱값 $\ATH(f)$는 Terhardt의 근사식으로 표현된다.
\begin{equation}\label{eq:ath}
  \ATH(f) = 3.64\left(\frac{f}{1000}\right)^{-0.8}
          - 6.5\, e^{-0.6\left(\frac{f}{1000} - 3.3\right)^{2}}
          + 10^{-3}\left(\frac{f}{1000}\right)^{4}
  \quad [\mathrm{dB\ SPL}].
\end{equation}
여기서 $\ATH(f)$는 주파수 $f$의 순음이 조용한 환경에서 겨우 들리는 음압
레벨(\dBSPL)이고, $f$는 \si{\hertz} 단위의 주파수이다. 세 항은 각각 저주파 상승,
\SIrange{2}{5}{\kilo\hertz} 부근의 민감도 향상(감소항), 고주파 급상승을 나타낸다.

\subsection{Bark(임계대역) 척도 변환 (Zwicker)}
선형 주파수 $f$를 임계대역 척도(Bark)로 변환하는 Zwicker의 근사식은 다음과 같다.
\begin{equation}\label{eq:bark}
  z(f) = 13\,\arctanop\!\bigl(0.00076\, f\bigr)
       + 3.5\,\arctanop\!\left(\left(\frac{f}{7500}\right)^{2}\right)
  \quad [\mathrm{Bark}].
\end{equation}
여기서 $z(f)$는 주파수 $f$에 대응하는 Bark 값이며, 가청 대역 전체는 약 24\,Bark 로
사상된다. 대응하는 임계대역폭(critical bandwidth) $\mathrm{BW}_{c}(f)$는 다음
근사식으로 주어진다.
\begin{equation}\label{eq:cbw}
  \mathrm{BW}_{c}(f) = 25 + 75\left[\,1 + 1.4\left(\frac{f}{1000}\right)^{2}\right]^{0.69}
  \quad [\si{\hertz}].
\end{equation}
$\mathrm{BW}_{c}(f)$는 주파수 $f$ 부근에서 하나의 임계대역이 차지하는 대역폭(\si{\hertz})
이다. 저주파에서는 약 \SI{100}{\hertz}로 거의 일정하다가 \SI{500}{\hertz} 이상에서 대략
주파수에 비례하여 증가한다.

\subsection{확산 함수 (Spreading Function)}
하나의 마스커가 이웃한 Bark 대역으로 미치는 마스킹 효과는 \emph{확산
함수}(spreading function)로 모형화한다. Bark 축 상의 거리
$\dz = z_{\text{maskee}} - z_{\text{masker}}$에 대한 Schroeder의 2-기울기(two-slope)
확산 함수는 다음과 같다.
\begin{equation}\label{eq:spread}
  \mathrm{SF}(\dz) = 15.81 + 7.5\,(\dz + 0.474)
                   - 17.5\,\sqrt{1 + (\dz + 0.474)^{2}}
  \quad [\si{\decibel}].
\end{equation}
여기서 $\dz$는 피마스커와 마스커의 Bark 위치 차이(\Bark)이며, $\mathrm{SF}(\dz)$는
그 상대 거리에서의 마스킹 감쇠(\si{\decibel})이다. $\dz>0$(고주파 방향)에서 완만히,
$\dz<0$(저주파 방향)에서 가파르게 감소하여 마스킹의 비대칭성을 나타낸다.

\subsection{톤성/비톤성 마스커 판별}
스펙트럼 성분 $S(k)$($k$는 이산 주파수 인덱스)가 국소 최대(local maximum)이고 이웃
성분보다 일정 기준 이상 강하면 \emph{톤성 마스커}(tonal masker)로 본다. MPEG-1 모델
1의 판별 기준은 다음과 같다.
\begin{equation}\label{eq:tonal}
  S(k)\ \text{is tonal} \iff
  \begin{cases}
    S(k) > S(k\pm 1), \ \text{및}\\[2pt]
    S(k) - S(k+j) \ge 7\dB \ \ \forall\, j \in \Delta_{k},
  \end{cases}
\end{equation}
여기서 $S(k)$는 인덱스 $k$에서의 음압 레벨(\dBSPL), $\Delta_{k}$는 검사 대상 이웃
인덱스 집합(주파수 대역에 따라 달라짐)이다. 톤성 성분을 이루는 인덱스들의 에너지를
합쳐 하나의 톤성 마스커 $P_{\mathrm{TM}}(k)$로 나타내고, 나머지(잔여) 에너지는 각
임계대역별로 모아 \emph{비톤성(잡음성) 마스커} $P_{\mathrm{NM}}(k)$로 취급한다.

\subsection{개별 마스킹 문턱값}
마스커 $j$(톤성 또는 비톤성)가 피마스커 위치 $i$에 만드는 개별 마스킹 문턱값
$T_{j}(i)$는 다음과 같이 표현된다.
\begin{equation}\label{eq:indiv}
  T_{j}(i) = P_{j}(z_{j}) + \mathrm{av}_{j}(z_{j}) + \mathrm{SF}\bigl(z_{i}-z_{j}\bigr),
\end{equation}
여기서 $P_{j}(z_{j})$는 Bark 위치 $z_{j}$에서 마스커의 음압 레벨(\dBSPL),
$z_{i}$는 피마스커의 Bark 위치, $\mathrm{SF}(\cdot)$는 \Cref{eq:spread}의 확산 함수,
$\mathrm{av}_{j}$는 마스킹 지수(masking index)이다. 마스킹 지수는 톤성/비톤성에 따라
다르며 대략
\begin{align}
  \mathrm{av}_{\mathrm{TM}} &= -1.525 - 0.275\,z - 4.5 \quad [\si{\decibel}], \label{eq:avtm}\\
  \mathrm{av}_{\mathrm{NM}} &= -1.525 - 0.175\,z - 0.5 \quad [\si{\decibel}], \label{eq:avnm}
\end{align}
로 주어진다($z$는 해당 마스커의 Bark 위치). 톤성 마스커가 비톤성 마스커보다 마스킹
능력이 약함(더 큰 음의 지수)을 반영한다.

\subsection{전역 마스킹 문턱값}
피마스커 위치 $i$에서의 \emph{전역 마스킹 문턱값}(global masking threshold)
$T_{g}(i)$는 절대 청취 문턱값과 모든 개별 마스킹 문턱값을 전력(에너지) 영역에서
합산하여 얻는다.
\begin{equation}\label{eq:global}
  T_{g}(i) = 10\log_{10}\!\left(
     10^{\,\ATH(i)/10}
     + \sum_{l=1}^{L} 10^{\,T_{\mathrm{TM},l}(i)/10}
     + \sum_{m=1}^{M} 10^{\,T_{\mathrm{NM},m}(i)/10}
  \right) \quad [\si{\decibel}].
\end{equation}
여기서 $\ATH(i)$는 \Cref{eq:ath}의 절대 청취 문턱값, $T_{\mathrm{TM},l}(i)$와
$T_{\mathrm{NM},m}(i)$는 각각 $L$개의 톤성, $M$개의 비톤성 마스커가 만드는 개별
문턱값(\Cref{eq:indiv})이다. 데시벨 값은 대수 척도이므로 반드시 선형 전력으로
환산하여 더한 뒤 다시 데시벨로 변환한다.

각 서브밴드 $n$에 대해 그 대역 내 최소 문턱값을 취하여
\emph{최소 마스킹 문턱값}(minimum masking threshold)을 정의한다.
\begin{equation}\label{eq:minmask}
  T_{\min}(n) = \min_{i \in \text{band } n} T_{g}(i) \quad [\si{\decibel}].
\end{equation}
보수적으로 대역 내 최솟값을 취함으로써 그 대역에서 어떤 성분도 가청되지 않도록
보장한다.

\subsection{SMR, MNR 및 SNR의 관계}
서브밴드 $n$에서의 신호 음압 레벨을 $L_{s}(n)$이라 하면 \emph{신호 대 마스크
비}(signal-to-mask ratio)는
\begin{equation}\label{eq:smr}
  \SMR(n) = L_{s}(n) - T_{\min}(n) \quad [\si{\decibel}]
\end{equation}
로 정의된다. 양자화기가 서브밴드 $n$에 $R(n)$비트를 배정하여 얻는 신호 대 잡음
비를 $\SNR(n)=\SNR\bigl(R(n)\bigr)$라 하면, \emph{마스크 대 잡음 비}(mask-to-noise
ratio)는
\begin{equation}\label{eq:mnr}
  \MNR(n) = \SNR(n) - \SMR(n) \quad [\si{\decibel}]
\end{equation}
이다. $\MNR(n) \ge 0$이면 서브밴드 $n$의 양자화 잡음이 마스킹 문턱값 아래에 있어
지각되지 않는다. 비트 할당의 목표는 주어진 비트 예산 안에서 모든 서브밴드의
최소 $\MNR$을 최대화(또는 모든 대역에서 $\MNR \ge 0$을 확보)하는 것이다.

%==============================================================================
\section{알고리즘 및 구현 기법}\label{sec:algorithm}
%==============================================================================
ISO/IEC 11172-3(MPEG-1 Audio)은 두 가지 심리음향 모델을 정의한다. 모델 1은 주로
Layer I/II에, 계산량이 큰 모델 2는 주로 Layer III(MP3)에 권장되며, AAC도 모델 2 계열
(불예측성 기반)을 사용한다.

\subsection{그림: ATH와 마스킹 곡선}
\Cref{fig:masking}은 절대 청취 문턱값(\Cref{eq:ath})과, \SI{1}{\kilo\hertz} 부근에
놓인 하나의 톤성 마스커가 만드는 마스킹 곡선을 개념적으로 보여준다. 두 곡선의 상위
포락선(upper envelope)이 사실상 전역 마스킹 문턱값이 되며, 그 아래의 신호 성분은
지각되지 않는다.

\begin{figure}[htbp]
  \centering
  \begin{tikzpicture}
    \begin{axis}[
        width=0.85\textwidth, height=6.5cm,
        xlabel={주파수 (\si{\hertz}, 로그 축)},
        ylabel={음압 레벨 (\dBSPL)},
        xmode=log, xmin=20, xmax=20000,
        ymin=-10, ymax=90,
        grid=both, legend pos=north east,
        every axis plot/.append style={very thick},
    ]
      % 절대 청취 문턱값 (Terhardt 근사, eq:ath)
      \addplot[blue, domain=20:20000, samples=200] {
        3.64*(x/1000)^(-0.8)
        - 6.5*exp(-0.6*(x/1000 - 3.3)^2)
        + 0.001*(x/1000)^4
      };
      \addlegendentry{절대 청취 문턱값 $\ATH(f)$}
      % 1 kHz 마스커 및 개념적 마스킹 곡선(비대칭 삼각형)
      \addplot[red, dashed] coordinates {
        (200,-5) (500,15) (800,45) (1000,60)
        (1250,40) (2000,20) (4000,2) (8000,-8)
      };
      \addlegendentry{마스킹 곡선 (\SI{1}{\kilo\hertz} 마스커)}
      \addplot[only marks, mark=*, black] coordinates {(1000,60)};
    \end{axis}
  \end{tikzpicture}
  \caption{절대 청취 문턱값과 \SI{1}{\kilo\hertz} 톤성 마스커의 개념적 마스킹 곡선.
           두 곡선의 상위 포락선 아래 신호는 지각되지 않는다.
           여기서 파란색 청취 문턱값 곡선은 Terhardt 근사식(\Cref{eq:ath})으로 계산한 것이며,
           빨간색 마스킹 곡선은 확산 함수로 계산한 것이 아니라 형태만 나타낸 개념적·예시용 곡선이다.}
  \label{fig:masking}
\end{figure}

\subsection{MPEG-1 심리음향 모델 1}
모델 1의 표준 처리 단계를 \Cref{alg:pam1}에 정리한다.

\begin{algorithm}[htbp]
  \caption{MPEG-1 심리음향 모델 1}
  \label{alg:pam1}
  \begin{algorithmic}[1]
    \State \textbf{입력:} PCM 프레임 $x[n]$ (예: 512 표본)
    \State \textbf{출력:} 서브밴드별 신호 대 마스크 비 $\SMR(n)$
    \State Hann 창을 적용하고 FFT 수행: $X(k) \gets \mathrm{FFT}\bigl(w[n]\,x[n]\bigr)$
    \State 스펙트럼을 정규화하여 음압 레벨로 변환: $P(k) \gets 90.302 + 10\log_{10}|X(k)|^{2}$ \Comment{\dBSPL}
    \State 각 서브밴드의 음압 레벨 $L_{s}(n)$ 계산 (대역 내 최대 성분 기준)
    \State 절대 청취 문턱값 $\ATH(k)$ 를 \Cref{eq:ath} 로 산출하여 포함
    \State 톤성 마스커 식별: \Cref{eq:tonal} 의 국소 최대·$7\dB$ 기준 적용
    \State 잔여 에너지를 임계대역별로 모아 비톤성(잡음) 마스커 결정
    \State \textbf{마스커 데시메이션(decimation):}
    \For{each 마스커 $j$}
      \If{$P_{j} < \ATH(z_{j})$}
        \State 마스커 $j$ 제거 \Comment{청취 문턱값 이하}
      \ElsIf{$0.5\Bark$ 이내에 더 강한 마스커 존재}
        \State 약한 마스커 제거
      \EndIf
    \EndFor
    \For{each 남은 마스커 $j$ 및 피마스커 위치 $i$}
      \State 개별 마스킹 문턱값 $T_{j}(i)$ 계산 (\Cref{eq:indiv,eq:spread})
    \EndFor
    \State 전역 마스킹 문턱값 $T_{g}(i)$ 합성 (\Cref{eq:global})
    \State 서브밴드별 최소 마스킹 문턱값 $T_{\min}(n)$ 계산 (\Cref{eq:minmask})
    \For{each 서브밴드 $n$}
      \State $\SMR(n) \gets L_{s}(n) - T_{\min}(n)$ \Comment{\Cref{eq:smr}}
    \EndFor
    \State \Return $\SMR(n)$
  \end{algorithmic}
\end{algorithm}

\subsection{MPEG-1 심리음향 모델 2}
모델 2는 모델 1과 달리 톤성 여부를 \emph{불예측성 측도}(unpredictability measure)로
연속적으로 추정한다. \Cref{alg:pam2}에 핵심 절차를 요약한다.

\begin{algorithm}[htbp]
  \caption{MPEG-1 심리음향 모델 2 (핵심)}
  \label{alg:pam2}
  \begin{algorithmic}[1]
    \State \textbf{입력:} PCM 프레임, 이전 두 프레임의 스펙트럼 이력
    \State \textbf{출력:} 임계대역(threshold-partition)별 $\SMR$
    \State 창(window) 적용 후 FFT 로 크기 $r(k)$ 와 위상 $\phi(k)$ 계산
    \State 예측값 $\hat r(k), \hat\phi(k)$ 를 직전 두 프레임으로 외삽
    \State 불예측성 측도 $c(k)$ 계산: 예측 오차가 클수록 잡음성(1에 근접)
    \State 에너지 $e(b)$ 와 가중 불예측성 $c(b)$ 를 threshold partition $b$ 로 집약
    \State 확산 함수로 $e(b), c(b)$ 를 인접 대역에 확산
    \State 톤성 지수(tonality index) $t(b) \gets -0.43\ln c(b) - 0.299$ 를 $[0,1]$ 로 제한
    \State 요구 SNR: $\mathrm{SNR}(b) \gets t(b)\cdot 18 + (1 - t(b))\cdot 6$ \Comment{\dB, 톤성일수록 큰 값}
    \State 전력비 $\mathrm{bc}(b) \gets 10^{-\mathrm{SNR}(b)/10}$ 로 대역별 문턱값 $\mathit{nb}(b) \gets e_{\text{spread}}(b)\cdot \mathrm{bc}(b)$
    \State 절대 청취 문턱값과 비교하여 하한 적용: $\mathit{thr}(b) \gets \max(\mathit{nb}(b), \ATH(b))$
    \State 사전 반향(pre-echo) 제어를 위해 직전 프레임 문턱값으로 상한 조정
    \State 각 부호화 대역의 $\SMR$ 계산 후 \Return
  \end{algorithmic}
\end{algorithm}

모델 2의 주요 차이점을 \Cref{tab:model12}에 정리한다.

\begin{table}[htbp]
  \centering
  \caption{MPEG-1 심리음향 모델 1과 모델 2의 비교}
  \label{tab:model12}
  \begin{tabular}{@{}lll@{}}
    \toprule
    항목 & 모델 1 & 모델 2 \\
    \midrule
    톤성 판별      & 국소 최대 + $7\dB$ 기준 & 불예측성 측도 $c(k)$ \\
    톤성 지표      & 이산(톤/비톤)            & 연속 톤성 지수 $t(b)$ \\
    시간 정보      & 미사용                   & 이전 프레임 이력 사용 \\
    주 적용 계층   & Layer I / II             & Layer III (MP3), AAC \\
    사전 반향 제어 & 제한적                   & 프레임 간 문턱값 제어 \\
    계산 복잡도    & 낮음                     & 높음 \\
    \bottomrule
  \end{tabular}
\end{table}

\subsection{임계대역 경계 예시}
\Cref{tab:cb}는 Bark 척도에 따른 임계대역 경계(중심/폭)의 대표적인 예이다.

\begin{table}[htbp]
  \centering
  \caption{임계대역(Bark) 경계의 대표 값 (근사)}
  \label{tab:cb}
  \begin{tabular}{@{}rrrr@{}}
    \toprule
    Bark $z$ & 하한 (\si{\hertz}) & 중심 (\si{\hertz}) & 대역폭 (\si{\hertz}) \\
    \midrule
    1  & 20   & 60   & 80  \\
    5  & 400  & 450  & 110 \\
    10 & 1080 & 1175 & 190 \\
    15 & 2320 & 2500 & 380 \\
    20 & 6400 & 7000 & 1300 \\
    24 & 15500 & 20500 & 3500 \\
    \bottomrule
  \end{tabular}
\end{table}

\subsection{비트 할당과의 결합}
심리음향 모델이 산출한 서브밴드별 $\SMR(n)$은 양자화·비트 할당 단계로 전달된다.
할당기는 각 서브밴드에 비트를 추가할 때 얻는 $\SNR(n)$을 이용하여
$\MNR(n)=\SNR(n)-\SMR(n)$(\Cref{eq:mnr})을 계산한다. 반복적으로 현재
$\MNR$이 가장 낮은(가장 가청 위험이 큰) 서브밴드에 비트를 추가하는 방식으로,
주어진 비트 예산 내에서 최소 $\MNR$을 최대화한다. 목표는 모든 대역에서
$\MNR(n)\ge 0$을 확보하여 양자화 잡음을 마스킹 문턱값 아래로 유지하는 것이다.
비트가 부족하면 지각적으로 가장 덜 중요한 대역부터 열화를 허용한다.

%==============================================================================
\section{맺음말}
%==============================================================================
심리음향 모델은 인간 청각의 마스킹 특성을 정량화하여, 지각적으로 무의미한 성분을
식별하고 그 아래로 양자화 잡음을 숨김으로써 고효율 오디오 압축을 가능하게 한다.
본 문서는 배경 이론에서 출발하여 표준 수식(\Cref{sec:math})과 MPEG-1 모델 1/2
알고리즘(\Cref{sec:algorithm})까지를 정리하였다. 실제 구현 시에는 대상 부호화기
(MP3, AAC 등)의 필터뱅크 구조와 프레임 길이에 맞추어 문턱값을 재사상하고, 사전
반향 억제와 블록 전환(block switching) 등 시간적 처리를 함께 고려해야 한다.

%==============================================================================
\begin{thebibliography}{9}
%==============================================================================
\bibitem{zwicker}
  E.~Zwicker and H.~Fastl,
  \emph{Psychoacoustics: Facts and Models},
  Springer.

\bibitem{painter}
  T.~Painter and A.~Spanias,
  ``Perceptual Coding of Digital Audio,''
  \emph{Proceedings of the IEEE}, vol.~88, no.~4, pp.~451--515, 2000.

\bibitem{iso11172}
  ISO/IEC 11172-3,
  \emph{Information technology -- Coding of moving pictures and associated audio
  for digital storage media at up to about \SI{1.5}{\mega\bit\per\second} --
  Part 3: Audio} (MPEG-1 Audio).

\bibitem{terhardt}
  E.~Terhardt,
  ``Calculating virtual pitch,''
  \emph{Hearing Research}, vol.~1, no.~2, pp.~155--182, 1979.

\bibitem{bosi}
  M.~Bosi and R.~E.~Goldberg,
  \emph{Introduction to Digital Audio Coding and Standards},
  Kluwer Academic Publishers.
\end{thebibliography}

\end{document}
